\chapter{Lassa Fever}
\index{Lassa Fever}

\section*{Definition and Overview}
Lassa fever is an acute viral haemorrhagic illness caused by the Lassa virus, a single-stranded RNA virus belonging to the family \textit{Arenaviridae}. First identified in 1969 in the town of Lassa in Borno State, Nigeria, the disease is endemic to parts of West Africa, including Nigeria, Sierra Leone, Liberia, and Guinea. The primary reservoir for the virus is the natal multimammate mouse (\textit{Mastomys natalensis}), a rodent ubiquitous in West African households. Transmission occurs primarily through direct contact with rodent excreta, consumption of contaminated food, or contact with soiled household items. Person-to-person transmission, particularly in healthcare settings with inadequate infection control (nosocomial spread), also occurs through direct contact with infected blood, bodily fluids, or mucous membranes. While approximately 80\% of infections are asymptomatic or mild, the remaining 20\% manifest as severe multi-system disease with significant morbidity and mortality.

\section*{Epidemiology}
Lassa fever is a major public health challenge in West Africa, with an estimated 100,000 to 300,000 infections and approximately 5,000 deaths occurring annually. The distribution of the disease is closely tied to the habitat of its rodent vector, \textit{Mastomys natalensis}. Sporadic outbreaks occur throughout the year, with a notable increase in cases during the dry season (December to April) due to agricultural activities and rodent movement into human dwellings. The overall case fatality rate is approximately 1\% to 2\%, but this rises to 15\% to 20\% among patients hospitalised with severe disease. Pregnant women, especially those in their third trimester, face an extremely high risk of maternal death (exceeding 30\%) and near-universal fetal loss. Lassa fever is occasionally diagnosed in travellers returning to non-endemic countries, highlighting the global relevance of the disease.

\section*{Aetiology and Pathophysiology}
The pathogenesis of Lassa fever involves viral entry, replication, and widespread immune suppression leading to vascular dysfunction.
\begin{itemize}
    \item \textbf{Viral structure:} Lassa virus is an enveloped, bisegmented RNA virus with a sandy appearance under electron microscopy due to host ribosomes inside the virion.
    \item \textbf{Host reservoir interaction:} Natal multimammate mice shed the virus in their urine and faeces throughout their lifetimes without showing signs of clinical illness.
    \item \textbf{Cellular entry:} The virus enters host cells by binding to the alpha-dystroglycan receptor, which is expressed on the surface of most tissues.
    \item \textbf{Target cells:} The primary targets are antigen-presenting cells (macrophages and dendritic cells), leading to impaired cytokine responses and suppressed immune activation.
    \item \textbf{Vascular dysfunction:} Unlike other haemorrhagic fevers, endothelial damage in Lassa fever is mediated by viral replication and functional impairment rather than direct cytolysis, leading to capillary leak, oedema, and cardiovascular collapse.
\end{itemize}

\section*{Clinical Features}
The incubation period ranges from 6 to 21 days, and the onset of symptoms is typically gradual and non-specific.
\begin{itemize}
    \item Gradual onset of fever, general malaise, headache, and joint pain.
    \item Severe sore throat with exudative pharyngitis is a classic diagnostic clue.
    \item Gastrointestinal symptoms, including nausea, vomiting, diarrhoea, and retrosternal or abdominal pain.
    \item Facial swelling, conjunctival injection, and mucosal bleeding in severe cases.
    \item Neurological manifestations, including tremors, encephalopathy, seizures, and cerebellar ataxia.
    \item Bilateral or unilateral sensorineural hearing loss, which develops in approximately 25\% of recovered patients.
\end{itemize}

\section*{Differential Diagnosis}
Given its non-specific early presentation, Lassa fever must be distinguished from several endemic febrile illnesses in West Africa.
\begin{itemize}
    \item \textbf{Malaria:} Plasmodium falciparum malaria is the most common cause of acute fever and must be ruled out immediately via blood smear or rapid diagnostic test.
    \item \textbf{Typhoid fever:} Presents with high fever, abdominal pain, and relative bradycardia; distinguished by blood cultures.
    \item \textbf{Yellow fever:} A flavivirus infection causing fever, jaundice, and haemorrhage; distinguished by serology or PCR.
    \item \textbf{Ebola and Marburg virus diseases:} Filovirus infections with rapid progression, severe gastrointestinal symptoms, and higher case fatality rates.
    \item \textbf{Sepsis:} Severe bacterial infection presenting with systemic inflammatory response syndrome and shock.
\end{itemize}

\section*{When to Seek Medical Care}
Patients residing in or travelling from endemic West African regions who develop sudden fever, sore throat, or bleeding should seek urgent medical evaluation at a facility equipped for high-consequence pathogens.

\textbf{Contact your local emergency services for:}
\begin{itemize}
    \item Severe bleeding from the nose, gums, or gastrointestinal tract
    \item Shortness of breath or persistent chest pain
    \item Confusion, seizures, or loss of consciousness
    \item Signs of shock, such as cold, clammy skin, and rapid heart rate
    \item Severe vomiting or inability to tolerate oral fluids
\end{itemize}

\section*{Diagnosis and Investigations}
Laboratory confirmation must be conducted in high-containment Biosafety Level 4 (BSL-4) facilities due to the infectious nature of the specimen.
\begin{itemize}
    \item \textbf{Reverse Transcription Polymerase Chain Reaction (RT-PCR):} The gold standard for detecting viral RNA in the acute phase (first 7 to 10 days of illness).
    \item \textbf{Enzyme-Linked Immunosorbent Assay (ELISA):} Used to detect Lassa-specific IgM and IgG antibodies, confirming diagnosis in later stages of the disease.
    \item \textbf{Viral isolation:} Performed in BSL-4 laboratories but rarely used for routine diagnosis due to safety risks.
    \item \textbf{Supportive laboratory investigations:} Full blood count showing leucopenia or leukocytosis, thrombocytopenia, and elevated liver enzymes (especially aspartate aminotransferase, which carries prognostic significance).
\end{itemize}

\section*{Management}
Management involves early antiviral therapy, supportive care, and strict isolation protocols to prevent nosocomial transmission.
\begin{itemize}
    \item \textbf{Antiviral therapy:} Intravenous ribavirin is the mainstay of treatment, highly effective if initiated within the first 6 days of symptom onset.
    \item \textbf{Fluid and electrolyte management:} Close monitoring of fluid balance to avoid fluid overload while maintaining renal perfusion.
    \item \textbf{Supportive care:} Antipyretics (avoiding NSAIDs to reduce bleeding risk), analgesics, and blood transfusions or clotting factors if severe haemorrhage occurs.
    \item \textbf{Infection prevention and control:} Strict adherence to contact and droplet precautions, isolating suspected patients in single rooms and using personal protective equipment (PPE).
    \item \textbf{Audiological monitoring:} Baseline and follow-up hearing assessments for patients recovering from the infection.
\end{itemize}

\section*{Complications}
\begin{itemize}
    \item Permanent or temporary sensorineural deafness, which occurs in a quarter of all survivors and is unresponsive to ribavirin.
    \item Severe pleural effusion, pericarditis, and congestive heart failure.
    \item Acute kidney injury and electrolyte disturbances.
    \item Maternal mortality and miscarriage, with a near 95\% mortality rate for the fetus in utero.
    \item Encephalopathy, seizures, and post-viral neurological deficits.
\end{itemize}

\section*{Prognosis and Outcomes}
The prognosis of Lassa fever depends heavily on the severity of presentation and the timing of ribavirin administration. Early initiation of treatment significantly reduces mortality to less than 5\% in hospitalised cases. Poor prognostic markers include high viral load, severe hepatic impairment (AST > 150 U/L), encephalopathy, and third-trimester pregnancy. Hearing impairment is a common and debilitating long-term sequela, with partial recovery occurring in only about half of affected individuals.

\section*{Prevention and Self-Care}
\begin{itemize}
    \item Prevent rodents from entering homes by sealing holes, cracks, and structural gaps.
    \item Store food in rodent-proof containers and avoid spreading food waste near human dwellings.
    \item Practice hand hygiene using soap and water or alcohol-based hand rubs regularly.
    \item Avoid consuming wild rodents or handling them in agricultural settings.
    \item Maintain strict barrier nursing and isolate febrile patients with travel history to endemic regions.
\end{itemize}

\section*{Sources and Further Reading}
The following reputable organisations provide detailed information and clinical guidance on Lassa fever:
\begin{itemize}
    \item World Health Organization. \textit{Lassa fever epidemiology, transmission, and prevention.} https://www.who.int/
    \item Centers for Disease Control and Prevention. \textit{Lassa fever clinical overview and rodent reservoir control.} https://www.cdc.gov/
    \item European Centre for Disease Prevention and Control. \textit{Factsheet on Lassa fever for health professionals.} https://www.ecdc.europa.eu/
    \item Healthdirect Australia. \textit{Haemorrhagic fevers and travel precautions.} https://www.healthdirect.gov.au/
    \item Gavi, the Vaccine Alliance. \textit{Lassa fever vaccines in clinical development.} https://www.gavi.org/
\end{itemize}

\clearpage
